% Page 004

\seite{4}

und Sefferweich, die sog.\ Sieben Gemeinden. Es\ob
werden aus demselben jährlich nicht unbedeutende\ob
Erträge aus der Lohe und dem Holze erzielt. Wie und\ob
wann die \enquote{Sieben Gemeinden} auf dem Bann\ohb
bezirke von Sefferweich und Malbergweich\ob
berechtigt wurden und gemeinsame Interessen\ob
zu vertreten hatten, ist nicht zu erfahren, hängt\ob
aber höchst wahrscheinlich mit dem alten Lehens\ohb
wesen zusammen.

\margmark{Die politische\\Gemeinde von\\Sefferweich.}%
Die Gemeinde Sefferweich besteht aus dem Dorfe\ob
Sefferweich, Waxbrunnen, Staffelstein links\ob
der Straße und Heinzenhof. Seit dem Jahre\ob
1853 sind die Kinder von Waxbrunnen\ob
\margmark{Waxbrunnen.}%
wegen des schlechten und weiten Weges aus dem\ob
Schulverbande von Sefferweich in den von Neiden\ohb
bach überwiesen worden. Die Frequenz der\ob
Schule ist eine sehr verschiedene. Im Jahre 1851\ob
betrug die Schülerzahl 56 dann sank dieselbe\ob
bis auf einige 40 und in den letzten 10 Jahren\ob
stieg dieselbe beständig, so daß heute die Schule\ob
von ca 60 Kindern besucht wird. In den Jahren\ob
von 1875 bis 1885 wanderten nämlich viele\ob
\margmark{Auswanderungen\\nach\\Amerika.}%
Bürger nach Amerika aus, woher die Ver\ohb
minderung der Einwohnerzahl wohl herrühren\ob
wird. Sefferweich gehört zur Pfarrei Seffern,
\pend
